\documentclass[11pt]{article}
\usepackage[margin=1in]{geometry}
\usepackage{hyperref}
\usepackage{enumitem}
\usepackage{booktabs}
\usepackage{xcolor}
\usepackage{amsmath}
\title{ProbOS --- Month 8 Pre-Plan\\Medical Device Domain (Part 1: Research Gate)}
\author{Nisong Monyimba}
\date{\today}

\begin{document}
\maketitle

\section*{Month goal}
Determine, via real research BEFORE any model code, whether medical
device battery safety can be genuinely validated.

\section*{The real, open question, stated up front}
Implantable device batteries typically use lithium-iodide or
lithium-silver-vanadium-oxide chemistry -- NOT the lithium-ion
chemistry \texttt{BatteryModel2Cell} already models. Unlike
automotive, medical device may require a genuinely new model.

\section*{Week 30 -- Research gate}
\begin{itemize}[leftmargin=*]
  \item Search for real, citable implantable device battery safety
    data, matching the pharma model's Day 1 discipline
  \item Confirm or reject the chemistry-transfer assumption with real
    sources
\end{itemize}

\textbf{Study materials for this week:}
\begin{itemize}[leftmargin=*]
  \item FDA guidance documents on implantable device battery safety
  \item ISO 14708 (active implantable medical devices)
  \item A survey paper on implantable-device battery electrochemistry
    (post-2015, IEEE/electrochemistry journals)
\end{itemize}

\section*{Week 31 -- Decision point}
\begin{itemize}[leftmargin=*]
  \item If real data found: design the correct chemistry-specific
    model
  \item If not found: document honestly, pivot to Month 10's research
    gate early
\end{itemize}

\section*{Exit criteria}
An honest, evidence-based answer -- not a model built on an assumed
chemistry transfer.

\end{document}
